\section{Configuração Experimental} \label{sec:experimentos}

Para avaliar o desempenho do FLASC, foram conduzidos experimentos utilizando o FLSC \cite{li2022federated} como linha de base.
O objetivo central da avaliação é qualificar o impacto de substituir o hiperparâmetro $N$ por $\rho$ (Veja a sessão ...).
Para garantir uma comparação equivalente entre os dois algoritmos,
variamos os hiperparâmetros $\rho$ e $N$ de forma correlacionada, conforme detalhado na Tabela \ref{tab:configs}.
Note que $N = 2$ significa que metade dos modelos serão selecionados.
Para reproduzir o equivalente com $\rho$,

\begin{table}[H]
\centering
\caption{Configuração dos hiperparâmetros de seleção para comparação.}
\label{tab:configs}
\begin{tabular}{lcc}
\hline
\textbf{Cenário} & \textbf{FLSC ($N$)} & \textbf{FLASC ($\rho$)} \\ \hline
Seleção Mínima   & 1                   & 0                       \\
Seleção Intermediária & 2                   & 0,5                     \\
Seleção Máxima (Inclusiva)   & 4                   & 1,0                     \\ \hline
\end{tabular}
\end{table}

Para os experimentos, utilizamos o conjunto de dados CIFAR-10,
que composto por 60.000 imagens coloridas de $32 \times 32$ pixels, distribuídas em 10 classes distintas.

A implementação do FLASC utiliza o framework Flower (\textit{v1.24.0}) \cite{beutel2020flower} com pytorch para orquestrar a simulação.
O código fonte completo, incluindo as estratégias de baseline e os scripts de reprodução,
está disponível publicamente em um repositório no GitHub\footnote{\url{https://github.com/nogueirapv/SimuAprendizadoFederado}}.

Quando ao particionamento dos dados, foi utilizado uma distribuição de Dirichlet com $\alpha = 0,5$
para refletir um cenário onde a distribuição de dados é desequilibrada.
Cada cliente divide sua partição local em 80\% para treinamento e 20\% para avaliação.

O modelo arquitetural adotado consiste em uma Rede Neural Convolucional (CNN) simples,
composta por duas camadas de convolução seguidas por camadas totalmente conectadas.
A primeira camada convolucional possui 6 filtros de tamanho $5 \times 5$,
enquanto a segunda possui 16 filtros de mesmo tamanho, ambas seguidas por operações de \textit{max-pooling} de $2 \times 2$.
As camadas densas subsequentes possuem 120, 84 e, finalmente, 10 neurônios na camada de saída, correspondendo às classes do CIFAR-10.
Foi utilizada a função de ativação ReLU em todas as camadas ocultas e o otimizador SGD com momento de $0,9$.
A função de perda adotada tanto para o treinamento local quanto para a avaliação de desempenho é a Entropia Cruzada (\textit{Cross-Entropy Loss}).

Configuramos o ambiente de simulação com 10 nós (clientes) participando de 10 rodadas de treinamento federado.
Em cada rodada, o servidor seleciona todos os clientes para o treinamento e metade deles para a etapa de avaliação.
Estabelecemos um conjunto de $K=4$ modelos globais para ambos os algoritmos.
Os parâmetros de treinamento local incluem 3 épocas, taxa de aprendizado de 0,001 e tamanho de lote de 32 amostras.

\subsection{Protocolo de Avaliação}

A avaliação ocorre ao final de cada rodada de treinamento.
Nesta etapa, o servidor transmite os $K$ modelos globais atualizados para os clientes selecionados.
Cada cliente avalia individualmente o desempenho de todos os $K$ modelos em seu conjunto de dados de teste local.
O cliente então identifica o modelo com o melhor desempenho (menor perda empírica) e reporta a acurácia e a perda deste modelo específico ao servidor.

O servidor agrega essas métricas calculando a média ponderada pelo número de amostras locais de cada cliente.
Este protocolo garante que a métrica final reflita a capacidade do sistema em oferecer, dentro do conjunto de clusters, pelo menos uma solução altamente especializada para cada distribuição de dados local encontrada na rede.
Dessa forma, avaliamos não apenas a convergência dos modelos, mas a eficácia da especialização promovida pelo agrupamento.
